% T33 source: Введение.tex
The effect of AI tools on software development productivity cannot be
characterized by a single universal speedup factor. In a controlled experiment
with GitHub Copilot, participants who completed the scoped task spent, on
average, 55.8\% less elapsed time on it than participants in the control group
\cite{peng2023impact}. By contrast, in a randomized study of experienced
developers working on tasks in familiar, mature open-source projects, access
to tools available in early 2025 increased adjusted active implementation time
by 19\% \cite{becker2025measuring}. A follow-up to that study reported likely
improvements in more recent agentic tools but also identified a deeper
methodological problem: developers selected tasks with the availability of AI
in mind and worked with multiple agents concurrently, so the duration of an
individual task no longer characterized the productivity of the overall
process \cite{metr2026uplift}. These findings concern different tasks, tools,
participants, and metrics and therefore do not constitute a direct
contradiction. Taken together, they show that the effectiveness of AI must be
assessed for an explicitly specified way of working, rather than for a model
or product outside its organizational context.

For development planning, this qualification is insufficient. Even a reliable
estimate of the duration of an individual AI-assisted task does not determine
the makespan of a set of interdependent tasks. The autonomous intervals of
multiple agents may overlap, but tasks compete for a limited number of streams,
follow dependencies, require integration, and periodically require input from
a single developer for task specification, clarification, review, or
acceptance.
Consequently, a local speedup may leave the overall makespan unchanged if it
does not reduce the active constraint, and an additional agent may fail to
create useful parallelism. Conversely, even an operating mode in which an
individual task takes longer than its manual baseline may reduce the makespan
of the full workload through overlap among autonomous phases. Thus, a distinct
problem of process organization and scheduling intervenes between an observed
task-level AI effect and the project's elapsed-time outcome.

This article considers one \emph{human--agent cell}: one developer completes a
prespecified AI-assisted set of software tasks using up to $P$ agent streams,
continuing each task until its result meets a fixed acceptance criterion. The no-AI
baseline is sequential completion of the same workload by one developer.
Interactive, delegated, and fully autonomous operating modes are
distinguished. In interactive mode, the human and the tool work jointly on a
task in a single stream; in delegated mode, task specification and acceptance
require the human, while autonomous agent phases may occur between them; in fully
autonomous mode, there is no mandatory human-attention phase. The article
focuses primarily on interactive and delegated tasks. Fully autonomous mode is
admitted as a limiting case with zero mandatory human work, but the discovery
of new tasks and automatic acceptance of results are not modeled.

For delegated tasks, a local retention rule is assumed: once started, a task
retains its assigned agent slot until its result is accepted. If the agent
requests the attention of a busy developer, it waits and is not assigned
another task, while the other agents continue working independently. A new
task may begin only in an available slot, and a dependent task may begin only
after all of its predecessors have been accepted. This rule describes a
specific process architecture rather than agent systems in general. Reuse of
a blocked slot, preemption, multiple developers, separate constrained CI/API
resources, dynamic revelation of the work graph, and interproject scheduling
are outside the scope of the formulation. The choice between fully manual and
AI-assisted execution for individual tasks is likewise treated as an external
decision; the model schedules the already selected AI-assisted task set.

The research question is: \emph{under what conditions does a fully specified
configuration in which one developer works with multiple AI agents reduce the
makespan of a fixed set of interdependent software tasks relative to sequential
work without AI, and what constraints determine the achievable speedup?} A
configuration specifies the number of agent streams and the operating mode of
each task; a substantively complete comparison must also hold fixed or
explicitly vary the decomposition, dependency graph, phase profiles,
acceptance criterion, and overhead functions. The objective is the makespan of
an accepted result. Quality and computational cost are included only through
their effects on the duration of review, correction, and task completion and
are not separate optimization objectives.

The central thesis of the article is that \textbf{achievable speedup is a
property of a fully specified human--agent process, not the product of ``AI
speed'' and a nominal number of agents}. The makespan may be constrained by the
aggregate workload of the streams, an indivisible task, a critical path
induced by technological precedence, the sequential human resource, or a
mixed resource chain created by phase ordering and the queue for developer
attention. Thus,
shortening a task that is not on the current critical path may still be useful
if it releases a scarce stream; increasing available capacity may reduce the
makespan until saturation, whereas moving to a configuration with a larger
configured $P$ may increase it if integration or attention-switching overhead
rises at the same time. Decomposition is similarly ambiguous: it removes
artificial serialization and enables parallel work, but adds boundaries for
task specification, review, and integration.

To answer the research question, a deterministic phase-based model M0--M4 is
proposed. The normalized duration $k_i^{(r)}$ pertains to the ``task---tool---mode''
triplet and compares the full cycle required to obtain an accepted result with
the baseline duration of the same task. It is a descriptive normalization, not
a causal estimate in the absence of a separate identification design. At
level M4, it is decomposed into the autonomous agent and human components
$a_i^{(r)}$ and $h_i^{(r)}$, and their internal order is specified by a
sequence of phases. The hierarchy begins with ideally divisible aggregate work
at M0 and then successively incorporates task indivisibility at M1, the
dependency graph at M2, cross-stream integration overhead at M3, and an exact
schedule of phases with a shared human resource and local blocking at M4. This
construction separates simple necessary constraints from the combinatorial
losses of a particular schedule.

The article provides four related results.

First, it defines an operational unit of comparison: not an AI tool in
isolation, but a complete ``tasks---modes---resources---acceptance criterion''
scenario. This makes it possible to distinguish an observed task duration from
a forecast, available capacity from configured parallelism, and an autonomous
agent phase from fully autonomous mode.

Second, for levels M0--M2, conditions are derived that relate normalized
durations, parallelism, and the work structure. At M0, speedup is equivalent to
$P>\bar{k}_w$. At M1--M2, this condition remains necessary, while the classical
list-scheduling bound yields a sufficient criterion in terms of the relative
length of the longest task or critical path. This explicitly distinguishes the
idealized speedup estimate, the structural lower bound, and the attainable
makespan.

Third, for M4, the structural lower bound
\[
  B_4=\max\{W_4/P,L_4,\gamma(P)H\},
\]
is introduced, and a speedup ceiling imposed by the sequential human resource
is proved. The bound combines the workload of the streams, the original
critical path, and total human work, but in general it is not equal to the
optimal makespan $T^*$. The resource-augmented phase graph provides an exact
representation of $T^*$ and shows how the service order of human-attention
phases and the order of tasks on the agent slots create mixed paths that are
absent from the original task-DAG\@. This distinguishes a lower bound, a
certified optimum, and the makespan $T(\pi)$ of a particular policy.

Fourth, a reproducible computational package with an exact solver and a fixed
deterministic policy tests the internal logic of the model in series
EXP-00--EXP-08. It reproduces the illustrative scenarios, provides
counterexamples to the universal attainability of $B_4$, separates the effect of
available capacity from that of increasing configuration overhead, examines
granularity and its interaction with $P$, tests scale robustness, and evaluates
whether the directions of the effects transfer across synthetic DAGs. These results constitute
computational verification and sensitivity analysis of the formal model, not
an empirical estimate for a specific team or AI tool.

The mathematical elements of this construction are individually well known:
indivisible jobs, the critical path, constrained resources, a shared service
resource, and resource-ordering arcs belong to scheduling theory. The claimed
contribution is not to reintroduce these objects but to map them systematically
to the complete cycle of a software task with a task- and mode-specific
coefficient $k$, separate agent and human-attention phases, repeated requests
for developer attention, agent-slot retention, integration overhead, and an
acceptance criterion for the result. Empirical calibration of this mapping
remains an open problem.

The article follows this logic. After the problem setting and compact notation,
it compares empirical estimates of AI productivity, scaling models, scheduling
theory, and research on human attention. It then introduces the M0--M4
hierarchy and the exact resource representation. The computational part first
specifies the method and rules of evidence and then examines the illustrative
scenario and series EXP-00--EXP-08. The discussion relates the formal and
computational results to practical use of the model; separate sections record
the limitations, threats to validity, and directions for future work. The full
glossary, detailed calculations, and additional results are provided in the
appendices.
